\documentclass[11pt]{article}
\usepackage[margin=1.1in]{geometry}
\usepackage{amsmath,amssymb}
\usepackage{booktabs}
\usepackage{graphicx}
\usepackage[round]{natbib}
\usepackage[hidelinks]{hyperref}
\usepackage{url}

\title{Mechanistic Hazard Ensembles for Geopolitical Forecasting:\\
An Auditable Alternative to Crowd and Language-Model Forecasters}
\author{Vivek Yadav\thanks{Independent researcher. Correspondence: vkn797029@gmail.com.
Code, data-fetch scripts, and a one-command replication package:
\url{https://github.com/vivek797029/Cassandra}.}}
\date{July 2026 — working draft v0.1}

\begin{document}
\maketitle

\begin{abstract}
Language-model forecasting systems produce fluent probability estimates whose
provenance cannot be inspected; human forecasting crowds are accurate but
slow to converge and equally opaque about mechanism. We study a third design
point: a fully mechanistic forecasting core in which every probability is
computed --- never generated --- from calibrated stochastic mechanisms, wrapped
in an audit layer that makes each forecast bit-for-bit reproducible. We ask
two questions. (1)~\emph{Coverage}: what fraction of real-world geopolitical
forecasting questions admits mechanistic treatment at all? Compiling 121
resolved conflict, oil, and macroeconomic questions from the Autocast corpus
under a precision-first compiler, we find 32 (26\%) compile to executable
mechanisms --- 30 to continuous-time hazard clocks and 2 to simulation path
predicates --- while 89 are structurally out of scope (institutional events,
country-specific series, count thresholds), a taxonomy we report in full.
(2)~\emph{Skill}: in a leakage-guarded walk-forward evaluation (126 forecasts
on 28 hazard questions, 2015--2022), maximum-likelihood hazard clocks with
partial pooling achieve a mean Brier score of 0.0812 against 0.0856 for the
contemporaneous forecasting crowd --- statistical parity (cluster-bootstrap
95\% CI on the difference $[-0.055, 0.054]$) --- and outperform the crowd in
the early, information-poor regime (Brier 0.087 vs.\ 0.116 at 10\% of question
lifetime) where strategic warning is most valuable. The model's uncertainty
is informative about its own failures: the width of its
plausibility-ball bands rank-correlates with realized squared error at
$\rho = 0.75$. We conclude that for the mechanizable quarter of questions, an
auditable clock matches an experienced human crowd at a fraction of the
latency, and argue that coverage taxonomies of this kind should accompany any
claim about ``AI forecasting.''
\end{abstract}

\section{Introduction}
Probabilistic forecasts of geopolitical and macroeconomic events inform
decisions with large downside risk. Two producer classes dominate current
practice. \emph{Crowds} --- prediction platforms and forecasting tournaments
--- are well calibrated in aggregate \citep{tetlock2015} but converge slowly,
and their estimates carry no inspectable causal structure. \emph{Language
models}, increasingly evaluated as forecasters \citep{zou2022}, emit
probabilities whose derivation is not auditable even in principle: the number
is a token, not a computation.

This paper evaluates a third design point, implemented in an open-source
system (\textsc{Cassandra}/ARGUS): a mechanistic Monte Carlo core in which
probabilities are \emph{counted} from simulated ensembles or computed in
closed form from calibrated hazard processes, and an interface layer that is
structurally incapable of inventing numbers --- a closed-grammar parser maps a
natural-language question either to an executable mechanism or to an explicit
refusal, and every answer embeds a manifest (parameter hash, seed, inputs)
sufficient to reproduce it exactly.

Two empirical questions determine whether this design point matters:

\textbf{Q1 (coverage).} What fraction of the questions people actually ask
admits mechanistic treatment? A method that answers only toy questions is a
curiosity regardless of its accuracy.

\textbf{Q2 (skill).} On the mechanizable fraction, does a calibrated
mechanism match human crowds under a protocol that excludes information
leakage?

Our contributions: (i)~a precision-first \emph{question compiler} and the
resulting coverage taxonomy over 121 real, resolved questions; (ii)~a
leakage-guarded walk-forward evaluation of partially pooled hazard clocks
against the contemporaneous crowd, with cluster-bootstrap inference and a
lifetime-resolved skill curve; (iii)~an ablation analysis showing which
components earn their keep; and (iv)~a replication package in which every
reported number regenerates from one command.

\section{Related work}
Judgmental forecasting tournaments established that trained crowds produce
calibrated geopolitical probabilities \citep{tetlock2015, mellers2014}.
Proper scoring rules for such evaluations are classical
\citep{brier1950, gneiting2007}. The Autocast benchmark \citep{zou2022}
introduced the resolved-question corpus with crowd trajectories we use, and
evaluated language models \emph{against} the crowd --- our protocol
deliberately mirrors theirs so results are comparable in kind. Survival and
point-process models of conflict have a long lineage in quantitative peace
research \citep{ward2010}; our contribution is not the hazard model itself but
its embedding in an end-to-end auditable pipeline plus the coverage question,
which the forecasting literature typically leaves implicit. Conformal and
robustness-based uncertainty wrappers \citep{vovk2005} motivate our
plausibility-ball bands.

\section{System context}
\textsc{Cassandra} is a 12-phase analytic pipeline: a regime-switching world
simulator with 16 named parameters $\theta$ (escalation, de-escalation, and
spillover channels linking conflict hazards to oil, inflation, growth, and
default cascades), trained by Brier minimization on retrospective questions;
an adversary that perturbs $\theta$ within a plausibility ball to produce
robust bands; and a conversational layer in which a closed-grammar parser
routes questions to scored engines and \emph{abstains} otherwise. The
serving layer enforces the audit contract: every response carries a manifest
ID resolvable to the exact $\theta$, seed, and inputs that produced it. The
present paper evaluates the epistemically load-bearing component --- the
mechanism family --- on real outcomes; the engineering controls (200+ tests,
red-team gates, contamination canaries, CI-enforced reproducibility) are
described in the repository.

\section{Question compilation and the coverage taxonomy}
\label{sec:compiler}
From the 6{,}532-question Autocast corpus we select resolved binary questions
in the system's mechanistic domains (armed conflict, oil, macroeconomics; 121
questions, 2015--2022, each with its full crowd trajectory). A
\emph{precision-first} compiler maps each question either to an executable
mechanism or to a refusal with a stable reason code; a question compiles only
when its semantics are unambiguously expressible over modeled state.

\begin{table}[h]\centering
\caption{Coverage taxonomy over 121 real questions (the compiler's verdicts).}
\label{tab:coverage}
\begin{tabular}{lrl}
\toprule
Verdict & $n$ & Meaning \\
\midrule
\textsc{hazard clock}      & 30 & first-event time in window (conflict) \\
\textsc{path predicate}    & 2  & threshold crossing of simulated series \\
\midrule
\textsc{institutional}     & 21 & announcements, legislation, approvals \\
\textsc{non-mechanistic}   & 48 & no modeled state variable \\
\textsc{count threshold}   & 4  & casualty counts (needs intensity model) \\
\textsc{country channel}   & 4  & national series outside global model \\
\textsc{unparseable window}& 8  & no confident deadline extraction \\
\textsc{sub-quarter/degenerate} & 4 & window below resolution \\
\bottomrule
\end{tabular}
\end{table}

Twenty-six percent of real questions are mechanizable today
(Table~\ref{tab:coverage}). We regard the taxonomy itself as a result:
claims about ``AI forecasting'' rarely state which questions the evaluated
system can represent at all.

\section{The hazard-clock family}
Each compiled conflict question --- ``will an event of class $c$ (airstrike,
confrontation, deployment, ceasefire, military action) occur in window
$[t_0,t_1]$?'' --- is modeled as a continuous-time hazard:
\[
P(\text{event}) = 1 - e^{-\lambda_c\,\Delta}, \qquad
\Delta = \max\!\bigl(0,\ t_1 - \max(t_{\mathrm{asof}}, t_0)\bigr)\ \text{years},
\]
the same mechanism family as the simulator's conflict channels, with the
exposure clipped to the \emph{remaining} window at forecast time. Rates are
estimated by maximum likelihood from all questions resolved strictly before
the forecast date: with observations $(\Delta_i, y_i)$,
$\mathcal{L}(\lambda) = \sum_i y_i \log(1 - e^{-\lambda\Delta_i})
- (1-y_i)\lambda\Delta_i$ is concave with monotone gradient, so the MLE is
found by bisection and clamped to $[0.01, 20]$ yr$^{-1}$. Sparse classes are
partially pooled toward the global rate through $\alpha{=}2$
pseudo-observations. Robust bands reuse the simulator's plausibility-ball
semantics, $\lambda \mapsto \lambda e^{\pm\rho}$ with $\rho = 0.10$.

\section{Evaluation protocol}
For every hazard question and every as-of date on a fixed grid of lifetime
fractions $\{0.1, 0.3, 0.5, 0.7, 0.9\}$, the model is refitted from scratch on
questions resolved strictly before that date and scored on the eventual
resolution. Guards: the contemporaneous crowd probability must exist at the
as-of date (else the pair is dropped for \emph{all} systems --- comparisons are
paired throughout); at least five resolved training questions must precede it.
This yields 126 forecast--outcome pairs over 28 questions. Baselines: the
crowd at the same instant, the historical base rate over resolved questions,
and uniform~0.5. Inference uses a cluster bootstrap resampling
\emph{questions} (records within a question share an outcome), $10^4$
resamples.

\section{Results}
\begin{table}[h]\centering
\caption{Mean Brier over 126 paired walk-forward forecasts (28 questions,
2015--2022). Lower is better.}
\label{tab:main}
\begin{tabular}{lc}
\toprule
System & Brier \\
\midrule
Hazard ensemble (this work) & \textbf{0.0812} \\
Crowd (contemporaneous)     & 0.0856 \\
Base rate                   & 0.1055 \\
Uniform 0.5                 & 0.2500 \\
\bottomrule
\end{tabular}
\end{table}

\textbf{Parity with the crowd.} The ensemble's aggregate Brier is nominally
better than the contemporaneous crowd's (skill $+5.1\%$), but the
cluster-bootstrap 95\% CI on the crowd$-$model difference is
$[-0.055,\ 0.054]$ ($P(\text{model better}) = 0.60$): at $n{=}28$ questions
the correct claim is \emph{statistical parity}, not superiority. Both clear
the base-rate and uniform baselines decisively.

\textbf{The early-warning regime.} Skill is not uniform over a question's
life (Table~\ref{tab:lifetime}). At 10\% of lifetime --- before the crowd has
absorbed much evidence --- the mechanism is clearly ahead (0.087 vs.\ 0.116);
by 90\% the crowd, which by then has read the news, is ahead (0.056 vs.\
0.072). For strategic warning, the early regime is the one that matters.

\begin{table}[h]\centering
\caption{Brier by lifetime fraction: the mechanism leads early; the crowd
overtakes as evidence accumulates.}
\label{tab:lifetime}
\begin{tabular}{lccccc}
\toprule
Fraction of lifetime & 0.1 & 0.3 & 0.5 & 0.7 & 0.9 \\
\midrule
Hazard ensemble & \textbf{0.0866} & 0.0856 & \textbf{0.0853} & 0.0782 & 0.0721 \\
Crowd           & 0.1158 & \textbf{0.0785} & 0.1108 & \textbf{0.0724} & \textbf{0.0563} \\
$n$             & 24 & 24 & 24 & 26 & 28 \\
\bottomrule
\end{tabular}
\end{table}

\textbf{Uncertainty that predicts its own errors.} The plausibility-ball band
width rank-correlates with realized squared error at Spearman
$\rho = 0.754$: wide bands flag the forecasts that miss. Calibration is
concentrated where the events live --- 120 of 126 forecasts fall in the
$[0,0.1)$ probability bin with an observed frequency of 0.083, i.e.\ slightly
\emph{under}confident in the rare-event regime.

\textbf{Ablations.} Removing partial pooling barely moves the aggregate
(0.0818); collapsing all event classes to a single global clock slightly
\emph{improves} it (0.0781). We report this against ourselves: at 30
questions, class-level granularity does not yet pay for its variance, and the
pre-registered primary (classes + pooling) is retained to avoid post-hoc
selection. The class structure is the component we expect to earn its keep
only as the corpus grows.

\section{Limitations}
The evaluable set is small (28 questions, 126 pairs); parity claims carry the
CI reported, and no significance is asserted. The two path-predicate
questions await historical state vintages (as-of oil prices) and are excluded
from the primary analysis. Event classification and window extraction are
rule-based; misparses are excluded by construction (precision over recall),
which biases coverage down but protects skill estimates. The crowd baseline
inherits each platform's participation patterns; early-lifetime crowd values
sit near platform priors, which is precisely the regime where a mechanistic
prior helps --- we report the full lifetime curve so readers can weight
regimes as their application demands. Finally, hazard clocks are the
\emph{simplest} member of the simulator's mechanism family; the evaluation
bounds the value of the family from below.

\section{Reproducibility and disclosure}
Every number in this paper regenerates from
\texttt{scripts/fetch\_real\_data.sh} followed by
\texttt{python scripts/run\_realeval.py}; the results artifact
(\texttt{output/realeval\_results.json}) contains every forecast record. The
system's serving layer attaches a reproducibility manifest to every answer,
and the repository's CI re-runs the full gate suite (unit, PostgreSQL,
red-team, load, container) on every push. This system and paper were
developed with substantial AI assistance (Anthropic Claude); all design
decisions, data selections, and claims were reviewed by the author, who bears
sole responsibility for them.

\begin{thebibliography}{9}\small
\bibitem[Brier(1950)]{brier1950} Brier, G.~W. (1950). Verification of
forecasts expressed in terms of probability. \emph{Monthly Weather Review},
78(1), 1--3.
\bibitem[Gneiting \& Raftery(2007)]{gneiting2007} Gneiting, T., \&
Raftery, A.~E. (2007). Strictly proper scoring rules, prediction, and
estimation. \emph{JASA}, 102(477), 359--378.
\bibitem[Mellers et al.(2014)]{mellers2014} Mellers, B., et al. (2014).
Psychological strategies for winning a geopolitical forecasting tournament.
\emph{Psychological Science}, 25(5), 1106--1115.
\bibitem[Tetlock \& Gardner(2015)]{tetlock2015} Tetlock, P.~E., \&
Gardner, D. (2015). \emph{Superforecasting: The Art and Science of
Prediction}. Crown.
\bibitem[Vovk et al.(2005)]{vovk2005} Vovk, V., Gammerman, A., \&
Shafer, G. (2005). \emph{Algorithmic Learning in a Random World}. Springer.
\bibitem[Ward et al.(2010)]{ward2010} Ward, M.~D., Greenhill, B.~D., \&
Bakke, K.~M. (2010). The perils of policy by p-value: Predicting civil
conflicts. \emph{Journal of Peace Research}, 47(4), 363--375.
\bibitem[Zou et al.(2022)]{zou2022} Zou, A., Xiao, T., Jia, R., et al.
(2022). Forecasting future world events with neural networks.
\emph{NeurIPS Datasets and Benchmarks} (arXiv:2206.15474).
\end{thebibliography}

\end{document}
